\documentclass[tikz,border=18pt]{standalone}
\usepackage{luatexja}
\usepackage{tikz}
\usetikzlibrary{shapes.geometric,arrows.meta,positioning,calc}

\definecolor{wallgray}{RGB}{120, 130, 140}
\definecolor{badred}{RGB}{190, 30, 30}
\definecolor{goodblue}{RGB}{20, 90, 160}
\definecolor{patchgreen}{RGB}{30, 140, 50}

\begin{document}
\begin{tikzpicture}[>=Stealth, font=\small]

  % -------------------------------------------------------------
  % 左側：脆弱性（セキュリティホール）を突くマルウェアの侵入
  % -------------------------------------------------------------
  \node[draw=badred, fill=red!6, rounded corners=8pt, line width=1.5pt, minimum width=8.8cm, minimum height=6.0cm, align=center] (left_box) at (-4.8, 2.0) {};
  
  \node[above, font=\large\bfseries, text=badred] at (-4.8, 4.3) {【危険】脆弱性を放置したシステム};

  % 壁（城壁）に穴がある
  \fill[wallgray!40] (-8.5, -0.2) rectangle (-7.0, 3.8);
  \fill[wallgray!40] (-2.6, -0.2) rectangle (-1.1, 3.8);
  
  \node[draw=badred, fill=red!20, dashed, very thick, minimum width=4.4cm, minimum height=1.6cm, align=center] (hole) at (-4.8, 1.8) {
    \textbf{\textcolor{badred}{脆弱性 (Security Hole)}}\\[2pt]
    \footnotesize プログラムの設計ミス・不具合
  };

  \node[draw=badred, fill=red!30, rounded corners=4pt, align=center] (malware) at (-4.8, 3.5) {
    \textbf{マルウェア / ランサムウェア}\\[2pt]
    \footnotesize (未修正の穴を狙って侵入！)
  };

  \draw[->, very thick, badred] (malware.south) -- (hole.north);
  \node[below, font=\scriptsize, text=badred] at (-4.8, 0.4) {
    【結果】データ暗号化・身代金要求・情報漏洩！
  };

  % -------------------------------------------------------------
  % 右側：セキュリティパッチによる修復と多層防御
  % -------------------------------------------------------------
  \node[draw=goodblue, fill=blue!6, rounded corners=8pt, line width=1.5pt, minimum width=8.8cm, minimum height=6.0cm, align=center] (right_box) at (4.8, 2.0) {};

  \node[above, font=\large\bfseries, text=goodblue] at (4.8, 4.3) {【安全】更新（パッチ適用）後の堅牢なシステム};

  % 壁が修復されている
  \fill[wallgray!40] (1.1, -0.2) rectangle (2.6, 3.8);
  \fill[patchgreen!40, draw=patchgreen, very thick] (2.6, 1.0) rectangle (7.0, 2.6);
  \fill[wallgray!40] (7.0, -0.2) rectangle (8.5, 3.8);

  \node[align=center] at (4.8, 1.8) {
    \textbf{\textcolor{patchgreen}{\Large 堅牢なパッチ (修正済)}}\\[2pt]
    \footnotesize OS・アプリの最新アップデート
  };

  \node[draw=badred, fill=red!15, rounded corners=4pt, align=center] (blocked_malware) at (4.8, 3.5) {
    \textbf{マルウェアの攻撃}\\[2pt]
    \footnotesize (侵入経路が完全に塞がれている)
  };

  \draw[->, very thick, badred] (blocked_malware.south) -- ++(0, -0.4);
  \node[font=\Large\bfseries, text=patchgreen] at (4.8, 2.8) {× 侵入阻止};

  \node[below, font=\scriptsize, text=goodblue] at (4.8, 0.4) {
    【結果】ゼロデイ攻撃リスクを最小化しデータを保護！
  };

\end{tikzpicture}
\end{document}
